\chapter{Introdução}
\label{cap:introducao}

% Descreva aqui o contexto geral do trabalho, apresentando o tema,
% a problemática e a justificativa da pesquisa.

% [Insira o texto da introdução aqui]
O avanço do conhecimento científico depende, em grande medida, da capacidade de organizar e comunicar resultados de forma clara e padronizada. Instituições de pesquisa como a \gls{EMBRAPA} desempenham papel fundamental nesse processo, ao gerar e difundir conhecimento aplicado a diversos setores, entre eles o da \gls{agropecuaria}. Este documento apresenta, de maneira genérica, a estrutura típica de um trabalho técnico-científico, com o objetivo de demonstrar a formatação adotada pelo template, cujo documento final é gerado em formato \gls{PDF}.

\section{Justificativa}
\label{sec:justificativa}

% Apresente a relevância e a motivação para a realização deste trabalho.

% [Insira a justificativa aqui]
A escolha do tema justifica-se por sua relevância teórica e prática, bem como pela necessidade de aprofundar a compreensão sobre o assunto. A realização deste estudo contribui para preencher lacunas identificadas na literatura e para fornecer subsídios a futuras investigações e aplicações.

\section{Objetivos}
\label{sec:objetivos}

% Apresente os objetivos geral e específicos do trabalho.

\subsection{Objetivo Geral}
\label{sec:objetivo-geral}

% [Insira o objetivo geral aqui]
Investigar de forma sistemática os fatores associados ao problema em estudo, com vistas a ampliar o conhecimento disponível sobre o tema.

\subsection{Objetivos Específicos}
\label{sec:objetivos-especificos}

% Liste os objetivos específicos usando alíneas:
%
% \begin{alineas}
% 	\item Objetivo específico 1;
% 	\item Objetivo específico 2;
% 	\item Objetivo específico 3.
% \end{alineas}

% [Insira os objetivos específicos aqui]
Para alcançar o objetivo geral, foram definidos os seguintes objetivos específicos:

\begin{alineas}
	\item revisar a literatura pertinente ao tema investigado;
	\item descrever os procedimentos metodológicos adotados;
	\item analisar os dados obtidos e discutir os resultados;
	\item apresentar conclusões e recomendações para trabalhos futuros.
\end{alineas}